\documentclass[11pt]{scrartcl}
\usepackage{evan}

\newtheorem{game}{Game}
\crefname{game}{Game}{Games}
\Crefname{game}{Game}{Games}
\setcounter{game}{-1}
\usepackage{multicol}
\ihead{Nine-Point Circle}
\ohead{Berkeley Math Circle}

\begin{asydef}
  // Old-style angle mark, matching the long-retired olympiad.asy version:
  // a single arc "pie slice" of radius t*markscalefactor centered at B.
  real markscalefactor = 0.03;
  path anglemark(pair A, pair B, pair C, real t = 8) {
    pair M = t * markscalefactor * unit(A - B) + B;
    pair N = t * markscalefactor * unit(C - B) + B;
    return arc(B, M, N) -- B -- cycle;
  }

  defaultpen(fontsize(11pt));
\end{asydef}


\begin{document}
\title{BMC Intermediate II: Triangle Centers}
\date{January 20, 2015}
\maketitle

\section{The Eyeballing Game}
\begin{multicols}{2}
\begin{game}
  Given a point $O$, the locus of points
  $P$ such that the length of $OP$ is $5$.
  \begin{center}
    \begin{asy}
      size(4cm);
      pair O = Drawing("O", origin, dir(225));
      pair P1 = Drawing("P_1", dir(40), dir(40));
      pair P2 = Drawing("P_2", dir(130), dir(130));
      pair P3 = Drawing("P_3", dir(-80), dir(-80));
      draw(O--P1);
      draw(O--P2);
      draw(O--P3);
    \end{asy}
  \end{center}
\end{game}
\begin{game}
  Given a triangle $ABC$, the locus of points $P$ such that $PB = PC$.
  \begin{center}
    \begin{asy}
      size(5cm);
      pair A = Drawing("A", dir(110), dir(110));
      pair B = Drawing("B", dir(210), dir(210));
      pair C = Drawing("C", dir(330), dir(330));
      pair P1 = Drawing("P_1", dir(90)*0.8, dir(75));
      pair P2 = Drawing("P_2", dir(90)*0.1, dir(45));
      pair P3 = Drawing("P_3", dir(90)*(-0.3), dir(-45));
      draw(B--P1--C, dotted);
      draw(A--B--C--cycle, heavyred);
    \end{asy}
  \end{center}
  \label{game:circumcenter}
\end{game}
\begin{game}
  Given a triangle $ABC$, the locus of points $P$ in $\triangle ABC$
  such that the distance from $P$ to $\ol{AB}$ and $\ol{AC}$
  is the same.
  \begin{center}
    \begin{asy}
      size(6cm);
      pair A = Drawing("A", dir(110), dir(110));
      pair B = Drawing("B", dir(210), dir(210));
      pair C = Drawing("C", dir(330), dir(330));
      draw(A--B--C--cycle, heavyred);
      pair P1 = Drawing("P_1", incenter(A,B,C), dir(-90));
      pair P2 = Drawing("P_2", extension(A,P1,B,C), dir(-90));
      draw(foot(P1,A,B)--P1--foot(P1,A,C), dashed);
      draw(foot(P2,A,B)--P2--foot(P2,A,C), dashed);
    \end{asy}
  \end{center}
  \label{game:incenter}
\end{game}
\columnbreak
\begin{game}
  Given a triangle $ABC$, the locus of points $P$
  such that the triangles $PAB$ and $PAC$ have the same area.
  (This one's a little harder, you might want to skip it
  and come back later.)
  \begin{center}
    \begin{asy}
      size(6cm);
      pair A = Drawing("A", dir(110), dir(110));
      pair B = Drawing("B", dir(210), dir(210));
      pair C = Drawing("C", dir(330), dir(330));
      draw(A--B--C--cycle, heavyred);
      pair P1 = Drawing("P_1", centroid(A,B,C), dir(-90));
      pair P2 = Drawing("P_2", midpoint(B--C), dir(-90));
      draw(A--P1--B, dashed);
      draw(A--P1--C, dashed);
    \end{asy}
  \end{center}
  \label{game:centroid}
\end{game}
%\begin{game}
%  Fix point $A$ on a circle $\omega$,
%  and vary a point $P$ on $\omega$.
%  What is the locus of midpoints $M$ of $AP$?
%  \begin{center}
%    \begin{asy}
%      size(6cm);
%      draw(unitcircle);
%      pair A = Drawing("A", (-1,0), dir(180));
%      pair P1 = Drawing("P_1", dir(80), dir(80));
%      pair P2 = Drawing("P_2", dir(10), dir(10));
%      pair P3 = Drawing("P_3", dir(-40), dir(-40));
%      pair M1 = Drawing("M_1", midpoint(Drawing(A--P1)), dir(140));
%      pair M2 = Drawing("M_2", midpoint(Drawing(A--P2)), dir(100));
%      pair M3 = Drawing("M_3", midpoint(Drawing(A--P3)), dir(100));
%    \end{asy}
%  \end{center}
%  \label{game:homotangent}
%\end{game}
\begin{game}
  Let $A$ be a point and consider a circle.
  Let $P$ be any point on the circle, and $M$ the midpoint of $\ol{AP}$.
  As $P$ varies, what is the locus of $M$?
  \begin{center}
    \begin{asy}
      size(6cm);
      draw(unitcircle, blue);
      pair A = Drawing("A", (-0.6,0), dir(140));
      pair P1 = Drawing("P_1", dir(80), dir(80));
      pair P2 = Drawing("P_2", dir(10), dir(10));
      pair P3 = Drawing("P_3", dir(230), dir(230));
      pair M1 = Drawing("M_1", midpoint(Drawing(A--P1)), dir(140));
      pair M2 = Drawing("M_2", midpoint(Drawing(A--P2)), dir(100));
      pair M3 = Drawing("M_3", midpoint(Drawing(A--P3)), dir(10));
    \end{asy}
  \end{center}
  \label{game:homocircle}
\end{game}
\end{multicols}
\eject

\section{Triangle Centers}
\begin{multicols}{2}
  \textbf{Circumcenter} $O$ (via \Cref{game:circumcenter})
  \begin{center}
    \begin{asy}
      size(6cm);
      pair A = Drawing("A", dir(110), dir(110));
      pair B = Drawing("B", dir(210), dir(210));
      pair C = Drawing("C", dir(330), dir(330));
      draw(A--B--C--cycle, heavyred);
      draw(unitcircle, blue);
      pair O = Drawing("O", origin, dir(45));
      draw(O--A); draw(O--B); draw(O--C);
      draw(O--(B+C)/2, dotted);
      draw(O--(C+A)/2, dotted);
      draw(O--(A+B)/2, dotted);
    \end{asy}
  \end{center}

  \textbf{Centroid} $G$ (via \Cref{game:centroid})
  \begin{center}
    \begin{asy}
      size(6cm);
      pair A = Drawing("A", dir(110), dir(110));
      pair B = Drawing("B", dir(210), dir(210));
      pair C = Drawing("C", dir(330), dir(330));
      draw(A--B--C--cycle, heavyred);
      pair G = Drawing("G", (A+B+C)/3, 1.4*dir(80));
      draw(A--Drawing(midpoint(B--C)), dashed);
      draw(B--Drawing(midpoint(C--A)), dashed);
      draw(C--Drawing(midpoint(A--B)), dashed);
    \end{asy}
  \end{center}

  \columnbreak

  \textbf{Incenter} $I$ (via \Cref{game:incenter})
  \begin{center}
    \begin{asy}
      size(6cm);
      pair A = Drawing("A", dir(110), dir(110));
      pair B = Drawing("B", dir(210), dir(210));
      pair C = Drawing("C", dir(330), dir(330));
      draw(A--B--C--cycle, heavyred);
      draw(incircle(A,B,C), blue);
      pair I = Drawing("I", incenter(A,B,C), dir(80));
      draw(A--I, dashed);
      draw(B--I, dashed);
      draw(C--I, dashed);
      draw(I--Drawing(foot(I,B,C)), dotted);
      draw(I--Drawing(foot(I,C,A)), dotted);
      draw(I--Drawing(foot(I,A,B)), dotted);
    \end{asy}
  \end{center}

  \textbf{Orthocenter} $H$ (via \Cref{game:circumcenter}! How?)
  \begin{center}
    \begin{asy}
      size(6cm);
      pair A = Drawing("A", dir(110), dir(110));
      pair B = Drawing("B", dir(210), dir(210));
      pair C = Drawing("C", dir(330), dir(330));
      draw(A--B--C--cycle, heavyred);
      draw(A--Drawing(foot(A,B,C)));
      draw(B--Drawing(foot(B,C,A)));
      draw(C--Drawing(foot(C,A,B)));
      pair H = Drawing("H", orthocenter(A,B,C), dir(70)*2);
    \end{asy}
  \end{center}
\end{multicols}

In words:
\begin{itemize}
  \ii Every triangle has exactly one circle passing through
  all of its vertices.
  This is called the \textbf{circumcircle},
  and its center $O$ is called the \textbf{circumcenter}.
  The point $O$ is the intersection of the perpendicular bisectors.

  \ii Every triangle has exactly one circle which is inscribed inside it.
  This is called the \textbf{incircle},
  and the center $I$ is called the \textbf{incenter}.
  It is the intersection of the angle bisectors.

  \ii The three \textbf{medians} of a triangle concur at
  a center $G$ called the \textbf{centroid} of the triangle.

  \ii The three \textbf{altitudes} of a triangle concur at
  a center $H$ called the \textbf{orthocenter} of the triangle.
\end{itemize}


\eject

\section{Orthocenter Flipping}
\begin{game}
  Given a triangle $ABC$, the points $P$ such
  that $\angle BPC = 180\dg - \angle A$.
  \begin{center}
    \begin{asy}
      size(6cm);
      pair A = Drawing("A", dir(110), dir(110));
      pair B = Drawing("B", dir(210), dir(210));
      pair C = Drawing("C", dir(330), dir(330));
      draw(A--B--C--cycle, heavyred);
      pair P1 = Drawing("P_1", orthocenter(A,B,C), dir(90));
      pair P2 = Drawing("P_2", -B*C/A, dir(-90));
      pair P3 = Drawing("P_3", -A, dir(-90));
      draw(B--P1--C, dashed);
      draw(B--P2--C, dashed);
      draw(B--P3--C, dashed);
    \end{asy}
  \end{center}
  \label{game:orthoflip}
\end{game}
Some hints for the above game:
\begin{itemize}
  \ii Can you identify the point $P_1$ I've drawn?
  \ii How are $P_2$ and $P_3$ related to $P_1$?
  \ii What's the relation between $A$, $B$, $P_2$, $P_3$, and $C$?
\end{itemize}

\section{Circle Theorems}
\begin{theorem}
  The sum of the angles in a triangle is $180\dg$.
\end{theorem}
\begin{theorem}
  [Inscribed Angle Theorem, also ``Star Trek Theorem'']
  An inscribed angle in a circle cuts out an arc equal to twice its measure.
\end{theorem}
\begin{center}
  \begin{asy}
    size(4cm);
    draw(unitcircle);
    pair A = dir(110);
    pair X = dir(210);
    pair Y = dir(330);
    pair B = dir( 50);
    Drawing("A", A, A);
    Drawing("B", B, B);
    Drawing("X", X, X);
    Drawing("Y", Y, Y);
    draw(X--A--Y);
    draw(X--B--Y);
    draw(anglemark(X,A,Y));
    draw(anglemark(X,B,Y));
  \end{asy}
  \quad
  \begin{asy}
    size(6cm);
    pair A = Drawing("A", dir(110), dir(110));
    pair B = Drawing("B", dir(210), dir(210));
    pair C = Drawing("C", dir(330), dir(330));
    draw(A--B--C--cycle);
    draw(unitcircle, blue);
    pair O = Drawing("O", origin, dir(50));
    draw(B--O--C);
    draw(O--A, dotted);
    markscalefactor *= 0.7;
    draw(anglemark(B,A,C));
    markscalefactor *= 0.7;
    draw(anglemark(B,O,C));
  \end{asy}
  \quad
  \begin{asy}
    size(4cm);
    draw(unitcircle);
    pair A = dir(110);
    pair X = dir(210);
    pair Y = dir(330);
    pair B = dir(280);
    Drawing("A", A, A);
    Drawing("B", B, B);
    Drawing("X", X, X);
    Drawing("Y", Y, Y);
    draw(X--A--Y);
    draw(X--B--Y);
    draw(anglemark(X,A,Y));
    pair Z = 1.5*B-0.5*X;
    draw(anglemark(Z,B,Y));
    draw(B--Z);
  \end{asy}
\end{center}

The inscribed angle theorem has important special cases.
As in the left diagram, suppose $A$, $B$, $Y$, $X$ lie on a circle in that order.
Then $\angle XAY = \angle XBY$, because both are equal to $\half \widehat{XY}$!
Similarly, if $A$, $X$, $B$, $Y$ lie on a circle in that order,
then $\angle XAY + \angle XBY = 180\dg$, because the two angles
mark out the entire circle: the sum of those arcs is $360\dg$.

It turns out the reverse is true: if points $A$, $B$, $X$, $Y$
satisfy the conditions above, all four lie on a circle.

\eject


\section{The Nine-Point Circle}
What happens when we combine \Cref{game:homocircle} and \Cref{game:orthoflip}?

\begin{theorem}
  [Nine-Point Circle]
  Let $ABC$ be a triangle with orthocenter $H$.
  The feet of the altitudes, the midpoints of the sides,
  as well as the three midpoints of $\ol{AH}$, $\ol{BH}$, $\ol{CH}$
  all lie on a single circle, called the \textbf{nine-point circle}
  of triangle $ABC$.
  Its center is the midpoint of $\ol{OH}$.
\end{theorem}

\begin{center}
  \begin{asy}
    size(12cm);
    pair A = Drawing("A", dir(110), dir(110));
    pair B = Drawing("B", dir(210), dir(210));
    pair C = Drawing("C", dir(330), dir(330));
    draw(A--B--C--cycle);
    draw(unitcircle, blue);
    pair H = Drawing("H", orthocenter(A,B,C), dir(45));
    pair A1 = Drawing("A_1", -B*C/A, dir(B+C));
    pair B1 = Drawing("B_1", -C*A/B, dir(C+A));
    pair C1 = Drawing("C_1", -A*B/C, dir(A+B));
    pair A2 = Drawing("A_2", -A, dir(B+C));
    pair B2 = Drawing("B_2", -B, dir(C+A));
    pair C2 = Drawing("C_2", -C, dir(A+B));

    pair M_A = Drawing("M_A", (B+C)/2, dir(B+C));
    pair M_B = Drawing("M_B", (C+A)/2, dir(C+A));
    pair M_C = Drawing("M_C", (A+B)/2, dir(A+B));
    pair D = Drawing("D", foot(A,B,C), dir(-A));
    pair E = Drawing("E", foot(B,C,A), dir(-B));
    pair F = Drawing("F", foot(C,A,B), dir(-C));
    pair H_A = Drawing(A+(B+C)/2);
    pair H_B = Drawing(B+(C+A)/2);
    pair H_C = Drawing(C+(A+B)/2);
    draw(circumcircle(D,E,F), dashed+red);
    draw(A--A1);
    draw(B--B1);
    draw(C--C1);
    draw(H--A2, dotted);
    draw(H--B2, dotted);
    draw(H--C2, dotted);

    pair O = Drawing("O", origin, dir(45));
    pair N9 = Drawing("N_9", midpoint(O--H), dir(65));
    draw(H--O, dotted);
  \end{asy}
\end{center}

Summary of work:
\begin{itemize}
  \ii $\angle BHC = 180\dg - \angle A$.
  \ii As we saw in \Cref{game:orthoflip}, the reflections of
  $H$ across $\ol{BC}$ and its midpoint all lie on the
  circumcircle of $\triangle ABC$.
  \ii Now apply \Cref{game:homocircle}.
\end{itemize}

\eject

\section{Euler's Line}
Now let $\overline{AM_A}$ intersect $\overline{OH}$ at $S$.

\begin{center}
  \begin{asy}
    size(12cm);
    pair A = Drawing("A", dir(110), dir(110));
    pair B = Drawing("B", dir(210), dir(210));
    pair C = Drawing("C", dir(330), dir(330));
    draw(A--B--C--cycle, blue);
    pair H = Drawing("H", orthocenter(A,B,C), dir(45));

    pair M_A = Drawing("M_A", (B+C)/2, dir(B+C));
    pair M_B = Drawing("M_B", (C+A)/2, dir(C+A));
    pair M_C = Drawing("M_C", (A+B)/2, dir(A+B));
    pair D = Drawing("D", foot(A,B,C), dir(-A));
    pair E = Drawing("E", foot(B,C,A), dir(-B));
    pair F = Drawing("F", foot(C,A,B), dir(-C));
    pair H_A = Drawing(A+(B+C)/2);
    pair H_B = Drawing(B+(C+A)/2);
    pair H_C = Drawing(C+(A+B)/2);
    draw(circumcircle(D,E,F), dashed+red);

    pair O = Drawing("O", origin, dir(45));
    draw(H--O, dotted+heavycyan);

    pair S = Drawing("S", centroid(A,B,C), dir(65));
    draw(A--H, heavycyan);
    draw(A--M_A, heavycyan);
    draw(O--M_A, heavycyan);
    draw(M_A--M_B--M_C--cycle, blue);
  \end{asy}
\end{center}

\begin{itemize}
  \ii Why are triangles $AHS$ and $OM_AS$ similar?
  \ii Why are triangles $ABC$ and $M_AM_BM_C$ similar? What ratio?
  \ii Which center of $M_AM_BM_C$ is $O$? (Not the circumcenter!)
  \ii Compute $AH / OM_A$.
  \ii Compute $HS / SO$ and $AS / SM_A$.
\end{itemize}

So let's give $S$ a new definition: it is the ``$\frac13$-way point of $O$ and $H$''. Our work above tells us that $A$, $S$, $M_A$ are collinear.
That means $BM_B$ and $CM_C$ must also pass through it.

So, what is $S$?

\begin{theorem}
  [Euler Line]
  The centroid, circumcenter, orthocenter, and nine-point center
  are all collinear.
\end{theorem}

\eject

\section{The Incenter}
% So we've got orthocenter, centroid, circumcenter

\begin{game}
  Let $ABC$ be a triangle with circumcircle $\Gamma$ and incenter $I$.
  We hold $\Gamma$, $B$, $C$ fixed and let $A$ vary on $\Gamma$.
  What is the locus of $I$?
  \begin{center}
    \begin{asy}
      size(7cm);
      pair B = Drawing("B", dir(210), dir(210));
      pair C = Drawing("C", dir(330), dir(330));
      draw(unitcircle, blue);
      draw(B--C, blue);

      pair A1 = Drawing("A_1", dir(130), dir(130));
      draw(B--A1--C);
      pair I1 = Drawing("I_1", incenter(A1,B,C), dir(A1));

      pair A2 = Drawing("A_2", dir(80), dir(80));
      draw(B--A2--C);
      pair I2 = Drawing("I_2", incenter(A2,B,C), dir(A2));

      pair A3 = Drawing("A_3", dir(40), dir(40));
      draw(B--A3--C);
      pair I3 = Drawing("I_3", incenter(A3,B,C), dir(A3));
    \end{asy}
  \end{center}
  Can you prove it?
\end{game}

To fully appreciate the result of this problem, I define the $A$-excenter $I_A$ to be the intersection of the \emph{external angles} of $\angle B$ and $\angle C$.
Then we get the following result.

\begin{theorem}
  [Incenter/Excenter Lemma]
  Let $ABC$ be a triangle with incenter $I$, $A$-excenter $I_A$, and denote by $L$ the midpoint of arc $BC$.  Then $L$ is the center of a circle through $I$, $I_A$, $B$, $C$.
\end{theorem}
\begin{center}
  \begin{asy}
    size(8cm);
    pair A = dir(110);
    pair B = dir(210);
    pair C = dir(330);
    pair I = incenter(A, B, C);
    draw(A--B--C--cycle);
    draw(unitcircle);
    pair L = dir(270);
    pair I_A = 2*L-I;
    draw(CP(L, I), blue);

    draw(A--I_A, dotted);

    dot("$A$", A, dir(A));
    dot("$B$", B, dir(190));
    dot("$C$", C, dir(-10));
    dot("$I$", I, dir(60));
    dot("$L$", L, dir(225));
    dot("$I_A$", I_A, dir(I_A));

    /* Source generated by TSQ */
  \end{asy}
\end{center}

You can do this entirely by just computing angles. Try it!

\eject

\section{Tying It All Together}
What happens if we add in all the excenters?

\begin{center}
  \begin{asy}
    size(12cm);
    pair A = dir(130);
    pair B = dir(210);
    pair C = dir(330);
    pair I = Drawing("I", incenter(A, B, C), dir(45));
    draw(A--B--C--cycle);
    draw(unitcircle, blue+dashed);
    Drawing("A", A, A);
    Drawing("B", B, B);
    Drawing("C", C, C);
    pair M_A = Drawing("M_A", dir(270), dir(225));
    pair M_B = Drawing("M_B", dir(50), dir(100));
    pair M_C = Drawing("M_C", dir(170), dir(200));
    pair I_A = Drawing("I_A", 2*M_A-I, dir(I-A));
    pair I_B = Drawing("I_B", 2*M_B-I, dir(I-B));
    pair I_C = Drawing("I_C", 2*M_C-I, dir(I-C));

    draw(I_A--I_B--I_C--cycle, heavyred);
    draw(A--I_A);
    draw(B--I_B);
    draw(C--I_C);
  \end{asy}
\end{center}



%%fakesection Not Using

%\section{A Couple Theorems I Won't Prove}
%\begin{multicols}{2}
%  \begin{theorem}
%    [Symmedian] In the figure below, $\angle BAX = \angle CAM$.
%  \end{theorem}
%  \begin{center}
%    \begin{asy}
%      pair A = Drawing("A", dir(130), dir(130));
%      pair B = Drawing("B", dir(210), dir(210));
%      pair C = Drawing("C", dir(330), dir(330));
%      draw(A--B--C--cycle);
%      draw(unitcircle);
%      pair M = Drawing("M", midpoint(B--C), dir(-90));
%      pair X = Drawing("X", 2*B*C/(B+C), dir(-90));
%      draw(B--X--C);
%      draw(X--A--M);
%      draw(anglemark(B,A,X));
%      draw(anglemark(M,A,C));
%    \end{asy}
%  \end{center}
%  \columnbreak
%  \begin{theorem}
%    [Pascal's Theorem] In the figure below, $X$, $Y$, $Z$ collinear.
%  \end{theorem}
%  \begin{center}
%    \begin{asy}
%      draw(unitcircle);
%      pair A,B,C,D,E,F;
%      A = Drawing(dir(100));
%      B = Drawing(dir(210));
%      C = Drawing(dir(30));
%      D = Drawing(dir(260));
%      E = Drawing(dir(140));
%      F = Drawing(dir(350));
%      draw(A--B--C--D--E--F--cycle);
%      pair X = Drawing("X", extension(A,B,D,E), dir(180));
%      pair Y = Drawing("Y", extension(B,C,E,F), dir(90));
%      pair Z = Drawing("Z", extension(C,D,F,A), dir(0));
%      draw(X--Z, dashed);
%    \end{asy}
%  \end{center}
%\end{multicols}

\section{A Problem To Think About If You're Bored}

\begin{center}
  \begin{asy}
    size(6cm);
    pair A = Drawing("A", dir(110), dir(110));
    pair B = Drawing("B", dir(210), dir(210));
    pair C = Drawing("C", dir(330), dir(330));
    pair M = Drawing("M", dir(-80), dir(-80));
    draw(A--B--C--cycle, heavyred+1);
    pair O = circumcenter(A,B,C);
    pair I = incenter(A,B,C);
    draw(unitcircle, blue);
    draw(incircle(A,B,C), blue+1);
    pair L = dir(90);
    pair T = Drawing("T", 2*foot(O,I,L)-L, dir(I-L));
    real pow = inradius(A,B,C)*inradius(A,B,C);
    pair Mc = I+pow/conj(M-I);
    pair[] K = intersectionpoints(incircle(A,B,C),
      Line(Mc-(M-I)*dir(90), Mc+(M-I)*dir(90)));
    pair K1 = K[1];
    pair K2 = K[0];
    pair X1 = Drawing("X_1", extension(K1,M,B,C), dir(-45));
    pair X2 = Drawing("X_2", extension(K2,M,B,C), dir(225));
    draw(K1--M--K2, darkgreen);
    pair MB = dir(40);
    pair MC = dir(160);
    pair TB = extension(T,MC,A,B);
    pair TC = extension(T,MB,A,C);
    draw(circumcircle(TB, TC, T), heavygreen);
    draw(circumcircle(T, X1, X2), cyan+dashed);
  \end{asy}
\end{center}
\begin{problem*}
  Let $ABC$ be a triangle and let $M$ be a point on its circumcircle.
  The tangents from $M$ to the incircle of $ABC$ intersect
  at points $X_1$ and $X_2$.
  A circle $\omega$ (called the $A$-mixtilinear incircle)
  touches $\ol{AB}$ and $\ol{AC}$, as well as the circumcircle internally
  at the point $T$.
  Prove that points $T$, $M$, $X_1$, $X_2$ lie on a circle.
\end{problem*}



Here is the diagram for the solution I gave. See if you can find some of the things we talked about hiding in this picture!
(For those of you that were here in October, there's also Pascal's Theorem
hiding somewhere in there.)
\begin{center}
  \begin{asy}
    size(16cm);
    pair A = Drawing("A", dir(110), dir(110));
    pair B = Drawing("B", dir(210), dir(210));
    pair C = Drawing("C", dir(330), dir(330));
    pair M = Drawing("M", dir(-80), dir(-80));
    draw(A--B--C--cycle, heavyred+1);
    pair O = circumcenter(A,B,C);
    pair I = incenter(A,B,C);
    draw(unitcircle, blue);
    draw(incircle(A,B,C), blue+1);
    pair L = Drawing("L", dir(90), dir(90));
    pair D = Drawing("D", foot(I,B,C), dir(B+C));
    pair E = Drawing("E", foot(I,C,A), dir(C+A));
    pair F = Drawing("F", foot(I,A,B), dir(A+B));
    pair H = Drawing("H", orthocenter(D,E,F), dir(160));
    draw(D--Drawing(foot(D,E,F)), darkmagenta);
    draw(D--E--F--cycle, blue+1);
    pair T = Drawing("T", 2*foot(O,I,L)-L, dir(I-L));
    pair Ac = Drawing("A^\ast", (E+F)/2, dir(60));
    pair Bc = Drawing("B^\ast", (F+D)/2, dir(B-I));
    pair Cc = Drawing("C^\ast", (D+E)/2, dir(0));
    draw(circumcircle(Ac,Bc,Cc), heavycyan+1);
    real pow = inradius(A,B,C)*inradius(A,B,C);
    pair Mc = Drawing("M^\ast", I+pow/conj(M-I), dir(100)*1.5);
    pair Lc = Drawing("L^\ast", I+pow/conj(L-I), dir( 10)*1.5);
    pair Tc = Drawing("T^\ast", I+pow/conj(T-I), dir(110)*1.5);
    Drawing("I", I, dir(45));
    draw(Ac--Bc--Cc--cycle, heavycyan);
    pair[] K = intersectionpoints(incircle(A,B,C),
      Line(Mc-(M-I)*dir(90), Mc+(M-I)*dir(90)));
    pair K1 = Drawing("K_1", K[1], dir(K[1]-K[0]));
    pair K2 = Drawing("K_2", K[0], dir(K[0]-K[1]));
    pair X1 = Drawing("X_1", extension(K1,M,B,C), dir(-45));
    pair X2 = Drawing("X_2", extension(K2,M,B,C), dir(225));
    pair X1c = Drawing("X_1^\ast", (D+K1)/2, dir(100));
    pair X2c = Drawing("X_2^\ast", (D+K2)/2, dir(80));
    draw(K1--D--K2, darkgreen);
    draw(K1--M--K2, darkgreen);
    draw(K1--K2, darkgreen);
    draw(Ac--Lc, darkmagenta);
    draw(L--T, green+dashed);
    draw(I--A, green+dashed);
    pair Dp = Drawing("D'", 2*Mc-D, dir(115));
    draw(D--Dp, dotted);
    pair Hp = Drawing("H'", 2*Mc-H, dir(-20));
    draw(H--Hp, dotted);
    pair MB = Drawing("M_B", dir(40), dir(40));
    pair MC = Drawing("M_C", dir(160), dir(160));
    pair TB = Drawing("T_B", extension(T,MC,A,B), dir(F-I));
    pair TC = Drawing("T_C", extension(T,MB,A,C), dir(E-I));
    draw(circumcircle(TB, TC, T), heavygreen);
    draw(T--MC--C, orange+dashed);
    draw(T--MB--B, orange+dashed);
    draw(MB--MC, heavymagenta+dashed);
    draw(A--L, heavymagenta+dashed);
    draw(TB--TC, orange+dotted);
    draw(circumcircle(T, X1, X2), cyan+dashed);
    draw(X2--I--X1, green+dotted);
    draw(I--M, green+dotted);
    draw( (MC+0.3*dir(B-A))--(MC+0.3*dir(A-B)), heavyred );
    draw( (MB+0.3*dir(C-A))--(MB+0.3*dir(A-C)), heavyred );
  \end{asy}
\end{center}

This problem is by Cosmin Poahata.
It appeared as the most difficult problem of the 2014 Taiwan Team Selection Tests for the International Mathematical Olympiad.
No one solved it during the time limit, although we discovered the above solution afterwards.

\end{document}
